% ---------------------------------------------------------------------------
% Project quality (45 points) -- Person A: structure and static quality
% (Maintainability: modularity, analysability, modifiability, testability).
% Person B's parts are marked "% >>> PERSON B" and are left as placeholders.
% Source of all numbers: personA/draft.md, raw output in personA/results/,
% reproducible with personA/measure/run.sh.
%
% Assignment brief (kept for reference, not typeset):
%   1. quality of the project as-is  2. concrete improvements
%   3. quality after the changes     4. trade-offs
%   Steps: (1) ISO 25010 selection (2) >=5 metrics, reproducible
%   (3) qualitative a.--e. with permalinks and UML (4) improvements traceable
%   to findings (5) re-assessment with before/after numbers.   Length: 3000 words.
% ---------------------------------------------------------------------------

\providecommand{\repocommit}{bf2323b}
\providecommand{\repocommitfull}{bf2323b3021c9d3db7a31e1e7d61068c3f704497}
\providecommand{\code}[1]{\texttt{\detokenize{#1}}}
\providecommand{\permalink}[2]{\href{https://github.com/anthropics/claude-code/blob/\repocommitfull/#1\##2}{\code{#1}}}
\providecommand{\todo}[1]{\textcolor{red}{[TODO: #1]}}

\section{Project quality (45 points)}
\label{sec:quality}

All measurements and file references in this section were taken at commit \texttt{\repocommit} (13 October 2026,
\href{https://github.com/anthropics/claude-code/tree/\repocommitfull}{permalink}). Line numbers from the interim report
were re-verified at this commit; where they had drifted they are corrected here.

%---------------------------------------------------------------------------
\subsection{Quality characteristics}
\label{sec:quality:characteristics}

% >>> PERSON B: Security, Compatibility, Reliability (fault tolerance), Flexibility (installability).

Under \textbf{Maintainability} we keep the interim report's \emph{modularity}, \emph{analysability} and \emph{testability}
and add \emph{modifiability}. The interim report argued that \code{mods/diff/} is modular because its code is spread
over hundreds of single-export files, but file count is evidence of modularity only if coupling is low, so we measure
it (M2). Modifiability is added because the DRY findings of the interim report and the concentration of change in one
file (37 of the 62 commits touching \code{mods/diff/hooks} since 9 September also touched \code{register.ts}) concern
the cost of change rather than structure. \emph{Reusability} is folded into modifiability: in this repository it appears
only as duplication (M3), and a separate heading would yield no separate proposal.

%---------------------------------------------------------------------------
\subsection{Metrics}
\label{sec:quality:metrics}

\paragraph{Scope.} We measure the TypeScript of \code{mods/{agents-md,diff,sec-default,telemetry}/hooks}, the tests in
\code{mods/*/tests}, \code{scripts/*.ts} and the Python of \code{plugins/}. We exclude \code{mods/types/claude-code.d.ts}
(13\,440 generated lines, which would dominate every number), Markdown plugin assets and JSON manifests.

% >>> PERSON B: M5--M7 (conformance, CI supply chain, fail-open rate)

\paragraph{M1: cyclomatic complexity and function length.}
\emph{Indicates} analysability, modifiability and testability: the number of paths a reader must understand and a test
must reach per unit. We used \texttt{lizard} for Python and, for TypeScript, \code{personA/measure/ccn.mjs}, a 60-line
McCabe counter on the TypeScript compiler AST. Lizard 1.24.1 ends a function early at a multi-line signature with a
generic return type (\texttt{): Promise<boolean> \{}); it silently drops \code{fetchBodies} in \code{register.ts}. CCN is
$1$ plus the number of \code{if}, \code{?:}, \code{case}, loops, \code{catch}, \code{&&}, \code{||} and \code{??}.

\begin{table}[H]
\centering
\small
\begin{tabular}{lrrrrrr}
\toprule
Corpus & Functions & Mean CCN & p90 & Max & CCN $>10$ & CCN $>15$ \\
\midrule
\code{mods/diff}       & 450 & 2.42 & 5  & 23 & 10 & 3  \\
\code{mods/sec-default}& 104 & 1.88 & 4  & 8  & 0  & 0  \\
\code{mods/telemetry}  & 125 & 2.43 & 5  & 16 & 2  & 1  \\
\code{mods/agents-md}  & 65  & 2.25 & 5  & 22 & 2  & 1  \\
\code{scripts/}        & 24  & 4.21 & 13 & 19 & 4  & 1  \\
\code{plugins/} (Py)   & 159 & 7.65 & 14 & 80 & 28 & 15 \\
\bottomrule
\end{tabular}
\caption{M1: function-level cyclomatic complexity per corpus at \texttt{\repocommit}.}
\label{tab:m1}
\end{table}

The worst functions are in \code{plugins/}: \code{handle_commit_review_posttooluse} (CCN 80, 476 lines,
\permalink{plugins/security-guidance/hooks/security_reminder_hook.py}{L901}), \code{handle_push_sweep_posttooluse} (62)
and \code{agentic_review} (61). In TypeScript the worst are \code{dockPlanOf} (23) and \code{autoCloseDuplicates}
(14, 174 lines, \permalink{scripts/auto-close-duplicates.ts}{L99}). The more important result is not in
\cref{tab:m1}: no function in \permalink{mods/diff/hooks/register.ts}{L1} exceeds CCN 14, yet the file defines 116
closures with \textbf{155 decision points}, against a median of 2 and a p90 of 9 for the other files of \code{mods/diff}
that contain functions, and is 3.4$\times$ the next most complex TypeScript file. Function-level CCN hides this because the
complexity is spread over closures sharing 25 pieces of mutable state, so we report the file-level figure next to it.
Complexity also splits along a governance boundary: \code{mods/} has mean CCN around 2.3, \code{plugins/} 7.65.

\paragraph{M2: coupling (fan-in / fan-out).}
\emph{Indicates} modularity and analysability. We ran \texttt{madge} 8.0.0 with \code{--ts-config mods/tsconfig.json}
over each mod's \code{hooks/} and computed fan-in and fan-out per file and per package with \code{personA/measure/fan.py}.
Every barrel file does \code{export * as default from '.'}, which madge reports as 267 ``circular dependencies''; these
are one-node self-loops and are removed first.

\begin{table}[H]
\centering
\small
\begin{tabular}{lrrrrr}
\toprule
Mod & Files & Edges & Mean fan-out & \code{register} fan-out & Multi-file cycles \\
\midrule
\code{diff}        & 594 & 1\,141 & 1.92 & 21 of 27 pkgs & 4 \\
\code{sec-default} & 83  & 130    & 1.57 & 8 of 9        & 0 \\
\code{telemetry}   & 210 & 365    & 1.74 & 6 of 16       & 1 \\
\code{agents-md}   & 56  & 103    & 1.84 & 6 of 7        & 0 \\
\bottomrule
\end{tabular}
\caption{M2: import coupling per mod (self-loops removed).}
\label{tab:m2}
\end{table}

The interim report's modularity claim survives for the leaves: mean fan-out is below 2, 11 of the 28 packages in
\code{diff} import nothing else, and the highest package fan-in is only 9 (\code{git}). A high fan-out for
\code{register} is not itself a defect, because a composition root imports everything. The defect is the combination
found only in \code{diff}: \code{register.ts} is a composition root, a stateful controller and a change hotspot at once
(1\,062 lines, 155 decision points = 24\,\% of the mod's total, 20 mutable \code{let}s). \code{telemetry} is the in-repo
counter-example: its \code{register} (6 decision points) only wires events to an object built by \code{telemetryOf}.
The four cycles in \code{diff} are \code{import type} edges erased at runtime, so they cost analysability, not correctness.

\paragraph{M3: duplication ratio.}
\emph{Indicates} modifiability. We used \texttt{jscpd} 5.4.1 with \code{--min-tokens 50} on TypeScript, TSX and Python.
\code{plugins/} has 2.98\,\% duplicated lines (11 clones), \code{scripts/} 4.94\,\% (3 clones) and \code{mods/} 4.02\,\%
(89 clones), but all 940 duplicated lines in \code{mods/} are repeated setup in test files; the production code of
\code{mods/} has \textbf{zero} clones, the strongest quantitative support for the interim report's DRY praise. The two
DRY violations found earlier are confirmed: the 27--28-line bootstrap block in four \code{hookify} entry scripts, and
\code{githubRequest<T>()}. jscpd finds only token-identical clones, so it misses the third copy of \code{githubRequest}
in \code{sweep.ts}, which has diverged (different signature, returns \code{{}} on a 404 where the others throw);
a divergent copy is worse than an identical one because a fix no longer applies mechanically.

\paragraph{M4: test distribution and effectiveness.}
\emph{Indicates} testability. We ran \code{claude plugin test mods/<mod>} (Claude Code 2.1.300, with
\code{CLAUDE_CODE_ENABLE_FUNCTION_HOOKS=1}). The runner has no coverage option, so we report structural proxies
instead of line coverage.

\begin{table}[H]
\centering
\small
\begin{tabular}{lrrrrr}
\toprule
Corpus & Source LOC & Test LOC & Test:source & Test cases & Pkgs with own test \\
\midrule
\code{mods/diff}        & 9\,512 & 7\,910 & 0.83 & 217 & 10 of 27 \\
\code{mods/sec-default} & 1\,320 & 4\,498 & 3.41 & 162 & 2 of 8 \\
\code{mods/telemetry}   & 3\,061 & 1\,952 & 0.64 & 34  & 1 of 14 \\
\code{mods/agents-md}   & 1\,071 & 370    & 0.35 & 5   & 0 of 6 \\
\code{plugins/} (Py, sh)& 9\,088 & 0      & 0    & 0   & 0 of 13 \\
\code{scripts/}         & 748    & 0      & 0    & 0   & --- \\
\bottomrule
\end{tabular}
\caption{M4: test volume and distribution at \texttt{\repocommit}. Test LOC includes fixtures.}
\label{tab:m4}
\end{table}

A package without its own test file is not necessarily untested, since each mod's \code{register.test.ts} drives the whole
mod through the engine; the last column counts only direct, isolatable tests. On Windows 11, \code{sec-default}
(162/162) and \code{agents-md} (5/5) pass, but \code{diff} failed 57/217 and then 31/217 in two identical runs, nearly
all by \code{timed out after 5000 ms}, plus three recurring assertion failures in \code{register}; \code{telemetry}
failed 8/34, including a Linux-only test. Failures that change between identical runs are environmental, so
\todo{replace with a Linux run (WSL or the CI run of \texttt{\repocommit}) before submission}. The metric also carries an
effectiveness caveat: \permalink{.github/workflows/mod-tests.yml}{L29-L32} exits 0 when the installed CLI lacks the test
command. With 2.1.300 the command exists, but a regression would turn the gate green without running anything, so a
pass rate is evidence only together with this caveat.

\paragraph{Reproducibility.}
\cref{tab:repro} gives tool, version and command for each metric; \code{personA/measure/run.sh} runs them all and
refuses to run unless \code{HEAD} is \texttt{\repocommit}.

\begin{table}[H]
\centering
\small
\begin{tabularx}{\textwidth}{llX}
\toprule
Metric & Tool (version) & Command \\
\midrule
M1 Python & lizard 1.24.1 & \code{lizard --csv plugins -l python} \\
M1 TypeScript & \code{ccn.mjs}, typescript 5.9, Node 25 & \code{node personA/measure/ccn.mjs mods/*/hooks scripts} \\
M2 & madge 8.0.0 & \code{madge --extensions ts --ts-config mods/tsconfig.json --json mods/<mod>/hooks} \\
M3 & jscpd 5.4.1 & \code{jscpd <dir> --format typescript,tsx,python --min-tokens 50} \\
M4 & Claude Code 2.1.300 & \code{claude plugin test mods/<mod>} \\
\bottomrule
\end{tabularx}
\caption{Tools and commands for M1--M4 (all at commit \texttt{\repocommit}, run on 13--19 October 2026).}
\label{tab:repro}
\end{table}

%---------------------------------------------------------------------------
\subsection{Qualitative assessment}
\label{sec:quality:qualitative}

\paragraph{(a) Architectural style.}
The interim report identified a microkernel (a closed core with extensions) and a tiered interceptor chain. The style suits a
product whose core is closed and whose extensions are developed in the open. The contract between them is one typed
declaration file, \code{mods/types/claude-code.d.ts}, against which every mod type-checks in CI
(\permalink{.github/workflows/mod-tests.yml}{L24-L25}). The style moves risk onto the extension boundary, and the repository
governs that boundary unevenly: \code{mods/} gets a typed contract, a test kit running each mod through the engine's real
hook chain, 418 test cases and no production duplication (M3), while \code{plugins/} gets a prose specification, no types and
no tests (M4), and holds the 15 functions with CCN $>15$ (M1). The style is well chosen, but its quality guarantee covers
only one of the two extension mechanisms the repository ships. The five-tier order is visible in the contract
(\code{TIERS = ["prepend","user","append","builtin","core"]}) and partly testable, since the test kit's \code{tier(...)}
seats a mod in a tier, but that the closed runtime \emph{enforces} the order cannot be verified from this repository.

% >>> PERSON B: (b) design principles, (c) role stereotypes, (e) documentation half

\paragraph{(d) Design patterns.}
\cref{tab:patterns} summarises the patterns carried over from the interim report. The central one is Model--View--Update (MVU),
applied correctly for three of its four roles. The \textbf{Model} is an immutable record
(\permalink{mods/diff/hooks/pane-state/pane-model/pane-model.ts}{L16}); \textbf{Update} is a set of pure reducers such as
\code{afterFetch}, each unit-tested without the engine; the \textbf{View}, \code{paneView}, is a pure function of the model.
The fourth role, the runtime that owns the model, sequences effects and feeds messages to Update, is not a component: it
is the body of \code{register(on)}, with 20 mutable \code{let}s plus 5 mutable containers, 10 hook registrations and 116
closures that reassign \code{model} directly in 22 places, 6 of them inside the View's own action callbacks
(\code{actionsOf}, L553--631), which mutate the model instead of dispatching a message. The pattern is three-quarters
applied, and the missing quarter is where M1 finds 155 decision points and where most recent changes land.
\cref{fig:lifecycle} shows that every step of a \code{/diff} lifecycle between the engine and the pure modules passes
through closure-local state.

\begin{table}[H]
\centering
\small
\begin{tabularx}{\textwidth}{llX}
\toprule
Pattern & Where & Verdict \\
\midrule
MVU & \code{mods/diff} & Model, Update, View clean; Controller absent (above) \\
Chain of Responsibility & tiered hook chain & Effective; order verifiable only at the contract \\
Decorator & hook wrappers in \code{mods/} & Effective, small and single-purpose \\
Strategy / Factory Method & \code{INSTALLED_BACKEND_PROBES} & Effective; a backend is added without editing the caller \\
Adapter & \code{$} $\rightarrow$ \code{Host} (21 members) & Effective; enables plain-fake unit tests \\
\bottomrule
\end{tabularx}
\caption{Design patterns and our judgement of their application.}
\label{tab:patterns}
\end{table}

\begin{figure}[H]
\centering
\IfFileExists{figures/diff-lifecycle-sequence.pdf}{\includegraphics[width=\textwidth]{figures/diff-lifecycle-sequence.pdf}}{\fbox{\todo{render personA/diagrams/diff-lifecycle-sequence.puml}}}
\caption{Sequence of one \code{/diff} lifecycle (session start, probe, fetch, reducer, view). Produced by manual static analysis of
\code{mods/diff/hooks/register.ts} and the packages it calls at \texttt{\repocommit}, rendered with PlantUML \todo{version}.}
\label{fig:lifecycle}
\end{figure}

\paragraph{(e) Test strategy.}
The suite mirrors the source (\code{tests/git/parse/parse-numstat.test.ts} $\leftrightarrow$ \code{hooks/git/parse/}), keeps
shared setup in \code{tests/fixtures/}, one export per file, and mixes two levels deliberately. Pure modules get unit tests
with no engine (\code{header-totals-of.test.ts}: 3 tests, 0 mocks). Each mod's \code{register.test.ts} is an integration test
through the real hook chain: every hook of the mod runs as shipped and anything the test leaves unanswered throws (in
\code{diff}, 45 such tests make 134 calls on \code{$.session/command/tool/ui}). This answers the interim report's open
question: the hook--engine contract \emph{is} tested, at integration level. The defects are elsewhere. \code{plugins/} (9\,088
LOC) and \code{scripts/} (748 LOC) have no tests (M4), and the CI gate can pass without running tests. In \code{diff}, 17 of
27 packages carry no direct test and depend on the integration tests, which are the slowest and most timing-sensitive; on our
machine every failure in three runs fell in 6 of the 30 test files.

%---------------------------------------------------------------------------
\subsection{Proposed improvements}
\label{sec:quality:improvements}

% >>> PERSON B: P4--P8

\paragraph{P1: complete MVU in \code{mods/diff} by decomposing \code{register.ts}.}
\emph{Motivated by} M1, M2, (a) and (d); \emph{targets} modularity, analysability, modifiability, testability.
\code{register.ts} becomes wiring only, as \code{telemetry}'s already is. A \code{PaneController} owns \code{model} and
dispatches each hook to one of five coordinators, each owning the closure variables that belong to it
(\cref{tab:p1}); View actions return messages to the controller instead of reassigning \code{model}
(\cref{fig:p1}). Trade-off: coordinators must read and write \code{model} across a boundary, and the
\code{isRefreshing} $\leftrightarrow$ \code{loadBodies} interleaving becomes a cross-object protocol that needs its own
test; the debounce and coalesce timers that keep \code{/diff} non-intrusive move intact into \code{RefreshScheduler}, so
their timing is unchanged.

\begin{table}[H]
\centering
\small
\begin{tabularx}{\textwidth}{lX}
\toprule
Coordinator & Closure state taken over (\code{register.ts} L43--71) \\
\midrule
\code{BodyLoader} (prototyped) & \code{bodyStamp}, \code{bodyBase}, \code{bodyLoads} \\
\code{RefreshScheduler} & \code{isRefreshing}, \code{isRefreshQueued}, \code{generation}, \code{landed}, \code{polled}, \code{timers} \\
\code{BackendPin} & \code{backend}, \code{probing}, \code{pin}, \code{loggedBaseKinds} \\
\code{PaneLifecycle} & \code{isPaneOpen}, \code{awaited}, \code{dialogRows}, \code{hasAutoOpened}, \code{hasRestoredEdits}, \code{columns}, \code{shownSessionId}, \code{sessionStartMs} \\
\code{AskComposer} & \code{armed}, \code{carrying} \\
\bottomrule
\end{tabularx}
\caption{P1: the five coordinators and the state each takes over.}
\label{tab:p1}
\end{table}

\begin{figure}[H]
\centering
\IfFileExists{figures/p1-after-component.pdf}{\includegraphics[width=\textwidth]{figures/p1-after-component.pdf}}{\fbox{\todo{render personA/diagrams/p1-after-component.puml}}}
\caption{Component view of \code{mods/diff} before and after P1; the shaded box is prototyped. Drawn manually from the
closure-variable assignments in \code{register.ts} at \texttt{\repocommit}, rendered with PlantUML \todo{version}.}
\label{fig:p1}
\end{figure}

\paragraph{P2: make \code{autoCloseDuplicates} testable.}
\emph{Motivated by} M1 and the interim report's testability finding; \emph{targets} testability and analysability. The
closure's decision logic is extracted into pure predicates (\code{isDupeDetectionComment}, \code{hasActivityAfter},
\code{hasAuthorDisagreed}, and two verdict functions). The comment-only verdict decides before reactions are fetched, so no
API call is added; the I/O loop keeps only fetching and acting.

\paragraph{P3: make test results trustworthy, then extend them to \code{plugins/}.}
\emph{Motivated by} M4 and (e); \emph{targets} testability. (1) Delete the \code{exit 0} branch in
\code{mod-tests.yml} L29--32 and pin the CLI version at L22, so a green gate implies the tests ran. (2) Raise the per-test
timeout or fix the timing-sensitive tests, since two identical runs differed by 26 failures. (3) Add \texttt{pytest}
suites for the two highest-risk plugin units: \code{hookify}'s rule engine (frontmatter extraction has CCN 32) and the
dispatch of \code{security_reminder_hook.py}, which needs a P2-style extraction first. (4) Add unit tests for the
\code{scripts/} predicates P2 creates. Trade-off: mandatory tests raise the bar for third-party plugin authors, a gain for
repository maintainability against a cost for ecosystem extensibility; we recommend requiring them for first-party
\code{plugins/} only.

%---------------------------------------------------------------------------
\subsection{Quality after the changes}
\label{sec:quality:reassessment}

P1 and P2 were applied to scratch copies, never to the repository, and measured with the commands of
\cref{tab:repro}. P3 and the B-owned proposals are discussed in the trade-off analysis below.

\begin{table}[H]
\centering
\small
\begin{tabular}{lrr}
\toprule
\code{scripts/auto-close-duplicates.ts} (P2) & Before & After \\
\midrule
Functions & 8 & 15 \\
Max CCN & 14 & 8 \\
Mean CCN & 3.62 & 2.47 \\
Functions with CCN $>10$ & 1 & 0 \\
Longest function (lines) & 174 & 49 \\
Total decision points & 21 & 22 \\
Unit tests possible without mocking the GitHub API & 0 & 6 (6 pass) \\
\bottomrule
\end{tabular}
\caption{P2 re-assessment, same counting rule as M1; the result type-checks under \code{tsc --strict}.}
\label{tab:p2}
\end{table}

\begin{table}[H]
\centering
\small
\begin{tabular}{lrr}
\toprule
\code{mods/diff/hooks/register.ts} (P1, \code{BodyLoader} only) & Before & After \\
\midrule
LOC (non-blank, lizard NLOC) & 735 & 701 \\
Decision points & 155 & 148 \\
Mutable \code{let} & 20 & 18 \\
Mutable containers & 5 & 4 \\
Package fan-out & 21 & 21 \\
\bottomrule
\end{tabular}
\caption{P1 re-assessment with one extracted coordinator (patch and module in \code{personA/prototypes/p1-body-loader/});
\code{tsc -p mods/tsconfig.json} passes.}
\label{tab:p1after}
\end{table}

P2 does not reduce total decision points (21 to 22); it relocates complexity into units that can be tested, which is the
honest claim: the longest function drops from 174 to 49 lines and the six new tests need no GitHub mock. For P1, extracting
one coordinator removes 7 decision points for 3 mutable items. \textbf{This is an estimate for the full refactor:} the
five coordinators take over 22 of the 25 mutable items, and at the same per-item ratio \code{register.ts} falls to roughly
100 decision points or fewer, a lower bound on the reduction since the hook bodies would also shrink to one-line
dispatches. Package fan-out should \emph{not} move: P1 improves state-per-unit, not coupling. As a behaviour check,
\code{claude plugin test diff} on the modified copy gave 165 pass and 52 fail, i.e.\ 49 timeouts plus the same three
assertion failures as both baseline runs and no new one. This is evidence, not proof, because a timed-out test never
reached its assertions; \todo{re-run on Linux before claiming unchanged behaviour}.

% >>> JOINT (S7): trade-off matrix of proposals x characteristics, covering P1 (performance, co-existence),
% >>> P3 (ecosystem extensibility) against B's P5 and P6.
